\documentclass[a4paper, serif, 10pt]{moderncv}

%% moderncv
% style and color
\moderncvstyle{banking}
\moderncvcolor{black}

% renew moderncv command to adapt for CJK colon
\renewcommand*{\cvitem}[3][.25em]{%
  \ifstrempty{#2}{}{\hintstyle{#2}：}{#3}%
  \par\addvspace{#1}}

\renewcommand*{\cvitemwithcomment}[4][.25em]{%
  \savebox{\cvitemwithcommentmainbox}{\ifstrempty{#2}{}{\hintstyle{#2}：}#3}%
  \setlength{\cvitemwithcommentmainlength}{\widthof{\usebox{\cvitemwithcommentmainbox}}}%
  \setlength{\cvitemwithcommentcommentlength}{\maincolumnwidth-\separatorcolumnwidth-\cvitemwithcommentmainlength}%
  \begin{minipage}[t]{\cvitemwithcommentmainlength}\usebox{\cvitemwithcommentmainbox}\end{minipage}%
  \hfill% fill of \separatorcolumnwidth
  \begin{minipage}[t]{\cvitemwithcommentcommentlength}\raggedleft\small\itshape#4\end{minipage}%
  \par\addvspace{#1}}

%% page layout/margins
\usepackage[top=2.5cm, bottom=2.5cm, left=1.5cm, right=1.5cm]{geometry}

\nopagenumbers{}

%% Babel config for English language
\usepackage[english]{babel}

%% fontspec
\usepackage{fontspec}

\IfFontExistsTF{Linux Libertine}{
  \setmainfont[Ligatures={TeX, Common}, Numbers=Lining, ItalicFont=Linux Libertine]{Linux Libertine}
}{}
\IfFontExistsTF{Linux Libertine O}{
  \setmainfont[Ligatures={TeX, Common}, Numbers=Lining, ItalicFont=Linux Libertine O]{Linux Libertine O}
}{}

%% CTeX
% CJK support, used to show CJK characters in the resume
%
% - fontset=none: disable builtin fontset but instead set the CJK font manually
% - heading=false: disable ctex heading
% - punct=kaiming: use kaiming punctuations styles for CJK
% - scheme=plain: use plain scheme, do not override \normalsize font size
% - space=auto: space settings for CJK characters
%
% ref:
% - http://ctan.mirrorcatalogs.com/language/chinese/ctex/ctex.pdf
\usepackage[UTF8, heading=false, punct=kaiming, scheme=plain, space=auto]{ctex}

\IfFontExistsTF{Noto Serif CJK SC}{
  \setCJKmainfont{Noto Serif CJK SC}
}{}
\IfFontExistsTF{Noto Sans CJK SC}{
  \setCJKsansfont{Noto Sans CJK SC}
}{}

%% line spacing
\usepackage{setspace}
\setstretch{1.125}

%% URL styling - use normal text instead of monospace
\urlstyle{same}

% Auto-underline all links
% Defer until \begin{document} so hyperref is already loaded
\AtBeginDocument{%
  \let\oldhref\href
  \renewcommand{\href}[2]{\oldhref{#1}{\underline{#2}}}%
}

%% Patch \httplink and \httpslink to handle full URLs
% The original moderncv macros always prepend http:// or https:// to URLs.
% This causes issues when the URL already has a protocol (e.g., https://example.com)
% resulting in malformed URLs like https://https://example.com.
% This patch checks if the URL already contains "://" and uses it directly if so.
% Note: We use \str_set:Nx (x-expansion) to expand macros like \@homepage first.
\ExplSyntaxOn
\str_new:N \g_yamlresume_url_str

\RenewDocumentCommand{\httpslink}{O{}m}{
  \str_set:Nx \g_yamlresume_url_str {#2}
  \str_if_in:NnTF \g_yamlresume_url_str {://}
    { \url{#2} }
    { \url{https://#2} }
}

\RenewDocumentCommand{\httplink}{O{}m}{
  \str_set:Nx \g_yamlresume_url_str {#2}
  \str_if_in:NnTF \g_yamlresume_url_str {://}
    { \url{#2} }
    { \url{http://#2} }
}
\ExplSyntaxOff

\name{佐藤 恵子}{}
\title{データサイエンティスト / 機械学習エンジニア}
\phone[mobile]{080-1234-5678}
\email{keiko.sato@example.jp}
\homepage{https://keiko-sato.dev}

\address{日本、東京都、東京、東京都千代田区丸の内1-2-3、100-0001}{}{}

\extrainfo{{\small \faGithub}\ \href{https://github.com/keikosato}{@keikosato} {} {} {} • {} {} {}
{\small \faLinkedin}\ \href{https://www.linkedin.com/in/keikosato}{@keiko-sato}}

\begin{document}

\maketitle

\section{基本情報}

\cvline{}{データサイエンティストとして、顧客分析・需要予測・レコメンドシステムの構築に5年以上従事。
ビジネス課題を統計分析・機械学習で解決し、施策の意思決定を支援しています。
%
\begin{itemize}
\item Python/Rを用いたデータ分析基盤の構築と運用
\item 機械学習モデルの設計・実装・評価（分類、回帰、クラスタリング、レコメンド）
\item ダッシュボード作成によるKPI可視化とインサイト共有
\item ビジネス部門と連携し、PoCから本番適用まで一気通貫で推進
\end{itemize}}
\section{学歴}

\cventry{2018年4月～2020年3月}
        {修士、統計学、成績：4.0}
        {東京大学大学院}
        {\url{https://www.u-tokyo.ac.jp/}}
        {}
        {\begin{itemize}
\item マーケティングデータのセグメンテーション手法に関する研究を行い、修士論文を執筆
\item 統計モデリングとPythonによる実データ分析のスキルを習得
\item 学会発表1回、論文投稿1件
\end{itemize}
\textbf{コース}：多変量解析、機械学習、ベイズ統計学、時系列解析、実験計画法、データ可視化}
\section{職歴}

\cventry{2021年4月～現在}
        {データサイエンティスト}
        {株式会社インサイトラボ}
        {\url{https://insightlab.example.com}}
        {}
        {\begin{itemize}
\item 小売業向け需要予測モデルを開発し、在庫削減と売上向上に貢献
\item 顧客セグメンテーションとLTV予測モデルを構築し、マーケティング施策のROI改善を実現
\item 各種分析ダッシュボードをLooker Studioで整備し、非エンジニアでもデータ活用できる環境を提供
\item データ基盤（BigQuery・dbt）の設計・運用を担当し、分析工数を50\%削減
\end{itemize}
\textbf{キーワード}：Python、BigQuery、Looker Studio、需要予測、セグメンテーション}

\cventry{2020年4月～2021年3月}
        {データアナリスト}
        {株式会社マーケットアナリティクス}
        {\url{https://market-analytics.example.com}}
        {}
        {\begin{itemize}
\item ウェブアクセスログと購買データを統合し、ユーザー行動分析を実施
\item アンケートデータの統計解析（クロス集計・回帰分析）による顧客満足度要因の特定
\item SQLやExcelを用いた定常レポートの自動作成により、レポート作成時間を短縮
\end{itemize}
\textbf{キーワード}：SQL、Excel、統計解析、レポーティング}
\section{言語}

\cvline{日本語}{ネイティブ・母国語レベル \hfill \textbf{キーワード}：母国語}
\cvline{英語}{完全な実務レベル \hfill \textbf{キーワード}：TOEIC 920、技術文書読解、プレゼン}
\section{スキル}

\cvline{データ分析・統計}{エキスパート \hfill \textbf{キーワード}：Python、R、SQL、統計モデリング、A/Bテスト}
\cvline{機械学習}{上級 \hfill \textbf{キーワード}：scikit-learn、LightGBM、TensorFlow、XGBoost、ハイパーパラメータ最適化}
\cvline{データ基盤・クラウド}{中級 \hfill \textbf{キーワード}：BigQuery、dbt、Airflow、Docker、Google Cloud}
\section{受賞・表彰}

\cventry{2022年12月}
        {株式会社インサイトラボ}
        {社内データ活用コンペ 最優秀賞}
        {}
        {}
        {小売需要予測の精度改善をテーマにした社内データコンペティションで最優秀賞を受賞。}
\section{資格・免状}

\cventry{2023年3月}
        {Amazon Web Services}
        {AWS Certified Machine Learning - Specialty}
        {\url{https://aws.amazon.com/certification/}}
        {}
        {}
\section{出版・著作}

\cventry{2020年9月}
        {マーケティングデータに基づく顧客セグメンテーション手法の比較検討}
        {データサイエンスジャーナル}
        {\url{https://example.com/publications/segmentation}}
        {}
        {実データを用いたクラスタリング手法の比較分析を行い、ビジネス適用における有効性を考察。}
\section{プロジェクト}

\cventry{2021年6月～2021年9月}
        {小売店舗向け需要予測および発注支援のためのWebダッシュボード}
        {需要予測ダッシュボード}
        {\url{https://github.com/keikosato/demand-forecast}}
        {}
        {\begin{itemize}
\item LightGBMとProphetによるアンサンブル予測
\item BigQueryでのデータ集計パイプライン構築
\item 店舗ごとの特性を考慮した特徴量エンジニアリング
\end{itemize}
\textbf{キーワード}：Python、LightGBM、Prophet、BigQuery、Dash}
\section{興味・関心}

\cvline{データ活用コミュニティ}{勉強会、ハッカソン、データ可視化}
\cvline{読書}{統計学、SF小説、技術書}
\section{ボランティア活動}

\cventry{2022年4月～現在}
        {運営スタッフ兼講師}
        {データ分析勉強会 Tokyo}
        {\url{https://example.com/study-group}}
        {}
        {\begin{itemize}
\item 初学者向けPythonデータ分析ハンズオンの講師を担当
\item 参加者のコードレビューやLT会の企画運営を通じて、コミュニティの活性化に貢献
\end{itemize}}

\end{document}